\documentclass{article}
\usepackage{geometry}
\geometry{a4paper,scale=0.8}
\usepackage[utf8]{inputenc}

\usepackage{amsmath, amssymb}
\usepackage{amsthm}
\usepackage{enumitem}
\usepackage{tikz-cd}
\usepackage{mathrsfs}

\newtheoremstyle{breakstyle}
    {5pt}
    {10pt}
    % {\normalfont}
    {\itshape}
    {0pt}
    {\bfseries}
    {\newline\indent}
    {0pt}
    {\thmname{#1}\thmnumber{ #2}\thmnote{ \textbf{[#3]}}}
\theoremstyle{breakstyle}
\newtheorem{exercise}{Exercise}[section]
\renewcommand{\theexercise}{\thesection.\Alph{exercise}}
\newtheorem{theorem}{Theorem}[section]
\newtheorem{proposition}{Proposition}[section]
\newtheorem{corollary}{Corollary}[section]
\newtheorem{remark}{Remark}[section]
\newtheorem{definition}{Definition}[section]
\newtheorem{lemma}{Lemma}[section]

\newenvironment{solution}
  {\begin{proof}[Solution.]\renewcommand{\qedsymbol}{}}
  {\end{proof}}

\DeclareMathOperator{\id}{Id}
\DeclareMathOperator{\mor}{Mor}
\DeclareMathOperator{\coker}{Coker}
\let\ker\relax
\DeclareMathOperator{\ker}{Ker}
\DeclareMathOperator{\im}{Im}
\let\lim\relax
\DeclareMathOperator{\lim}{Lim}
\DeclareMathOperator{\colim}{Colim}
\let\hom\relax
\DeclareMathOperator{\hom}{Hom}
\let\mod\relax
\DeclareMathOperator{\mod}{Mod}
\DeclareMathOperator{\pr}{pr}
\DeclareMathOperator{\bil}{Bil}
\let\dim\relax
\DeclareMathOperator{\dim}{Dim}
\let\gcd\relax
\DeclareMathOperator{\gcd}{Gcd}

\title{Chapter 1 - Just enough category theory to be dangerous}
\author{Flandre}
\date{\today}

\begin{document}
\maketitle

\section{Categories and functors}
\begin{exercise}
UNIMPORTANT EXERCISE.

A category in which each morphism is an isomorphism is called a groupoid.
(This notion is not important in what we will discuss. The point of this
exercise is to give you some practice with categories, by relating them to
an object you know well.)

\begin{enumerate}[label=({\alph*})]
\item A perverse definition of a group is: a groupoid with one object. Make
    sense of this. (Similarly, in case you care: a perverse definition of a
    monoid is: a category with one object.)
\item Describe a groupoid that is not a group.
\end{enumerate}

\begin{solution}
	\begin{enumerate}
        \item Since there's only one object, namely $e$, in the category, all
            the morphisms in such category will send $e$ to $e$.
			\begin{itemize}
                \item The combinations of morphisms are always associative.
                \item There must exist an identity morphism, which'll be the
                    identity in the group.
				\item Since it's a groupoid, every morphism has a inverse.
			\end{itemize}
            \emph{ (In fact every group can be represented using the language of
            vategory. Consider Cayley's Theorem and we ca see every group
            element can be regarded as an automorphism, satisfying the case of
            morphism to a object itself.) }
        \item We consider a category $\mathscr{C}$ whose objects are $A,B$.
            Consider
			$$
			\mor(\mathscr{C})=\{\id_{A},\id_{B}\}
			.$$ 
            and we can see that it's a groupiod where the elements
            corresponding to $\id_{A},\id_{B}$ cannot have multiplication ,
            therefore not a group.
	\end{enumerate}
\end{solution}

\end{exercise}

\begin{exercise}
If $A$ is an object in a category $\mathcal{C}$, show that the invertible
elements of $\mor(A,A)$ form a group (called the automorphism group of $A$,
denoted $\operatorname{Aut}(A)$). What are the automorphism groups of the
objects in Examples 1.1.2 and 1.1.3? Show that two isomorphic objects have
isomorphic automorphism groups. (For readers with a topological background:
if $X$ is a topological space, then the fundamental groupoid is the category
where the objects are points of $X$, and the morphisms $x\to y$ are paths
from $x$ to $y$, up to homotopy. Then the automorphism group of $x_{0}$ is
the (pointed) fundamental group $\pi_{1}(X,x_{0})$. In the case where $X$ is
connected, and $\pi_{1}(X)$ is not abelian, this illustrates the fact that
for a connected groupoid --- whose definition you can guess --- the
automorphism groups of the objects are all isomorphic, but not canonically
isomorphic.)

\begin{proof}
		\begin{enumerate}
            \item Similar to the case in \textbf{1.1.A}, we can use exactly the
                same proof to show that invertible elements of
                $\mor(A,A)$ forms a group.
            \item Automorphism groups of objects in \textbf{Example 1.1.2} are
                those permutation groups permutating the elements in a group.
            \item Automorphism groups of objects in \textbf{Example 1.1.3} are
                those invertible linear transformations mapping a linear space
                to itself.
		\end{enumerate}
\end{proof}

\end{exercise}

\begin{exercise}[Skipped, require \textbf{1.1.21} to solve.]
Let $(\cdot)^{\vee\vee}:f.d.Vec_{k}\to f.d.Vec_{k}$ be the double dual
functor from the category of finite-dimensional vector spaces over $k$ to
itself. Show that $(\cdot)^{\vee\vee}$ is naturally isomorphic to the
identity functor on $f.d.Vec_{k}$. (Without the finite-dimensionality
hypothesis, we only get a natural transformation of functors from id to
$(\cdot)^{\vee\vee}$.)

\begin{proof}
    \begin{enumerate}[label=({\roman*})]
        \item In the case of finite cases, the set of bases of
            $Vec_{k},Vec_{k}^{\vee}$ is finite. Therefore the functor
            $(\cdot)^{\vee}$ must be isomorphism.
	\end{enumerate}
\end{proof}

\end{exercise}

\begin{exercise}
Show that $\mathcal{V} \rightarrow f.d.Vec_{k}$ gives an equivalence of
categories, by describing an ``inverse'' functor. (Recall that we are being
cavalier about set-theoretic assumptions, see Caution 0.3.1, so feel free to
simultaneously choose bases for each vector space in $f.d.Vec_{k}$. To make
this precise, you will need to use G\"{o}del-Bernays set theory or else
replace $f.d.Vec_{k}$ with a very similar small category, but we won't worry
about this.)

\begin{proof}
    Since it's finite, just find out the bases and we can define a trivial
    isomorphism between $\mathscr{V}$ and $f.d.Vec_{k}$ by mapping those bases
    respectively.
\end{proof}

\end{exercise}

\section{Universal properties determine an object up to unique isomorphism}
\begin{exercise}
Show that any two initial objects are uniquely isomorphic. Show that any two
final objects are uniquely isomorphic.

\begin{proof}
	\begin{enumerate}
		\item For any two initial objects $A,B$, we define
			 $$
			\varphi: A\mapsto B,\psi: B\mapsto A
			.$$ 
			And we can see 
			$$
			\varphi \circ \psi :B\mapsto B,\psi\circ\varphi:A\mapsto A
			.$$ 
            Also, notice since $A,B$ are initial objects, they have only one
            unique map to every objects in the category, including themselves.
            And $Id_{A},Id_{B}$ maps $A,B$ to themselves. Hence we obtain
			\begin{align*}
				& Id_{A}=\psi\circ\varphi \\
				& Id_{B}=\psi\circ\varphi
			.\end{align*}
			So we know $\varphi,\phi$ are isomorphisms.
		\item The proof is similar to \emph{(a)}, skipped.
	\end{enumerate}
\end{proof}

\end{exercise}

\begin{exercise}
What are the initial and final objects in Sets, Rings, and Top (if they
exist)? How about in the two examples of \S\emph{1.1.9}?

\begin{solution}
	\begin{enumerate}
        \item In the category of sets, the initial and final object is
            $\emptyset$ and the class of all sets, respectively.
        \item In the category of rings, the initial and final object is the zero
            ring and the class of all rings, respectively.
		\item \emph{Skipped (haven't learn topology yet).}
	\end{enumerate}
\end{solution}

\end{exercise}

\begin{exercise}
Show that $A \to S^{-1}A$ is injective if and only if $S$ contains no
zerodivisors. (A zerodivisor of a ring $A$ is an element $a$ such that there
is a nonzero element $b$ with $ab = 0$. The other elements of $A$ are called
non-zerodivisors. For example, an invertible element is never a zerodivisor.
Counter-intuitively, $0$ is a zerodivisor in every ring but the $0$-ring.
More generally, if $M$ is an $A$-module, then $a \in A$ is a zerodivisor for
$M$ if there is a nonzero $m \in M$ with $am = 0$. The other elements of $A$
are called non-zerodivisors for $M$. Equivalently, and very usefully,
$a \in A$ is a non-zerodivisor for $M$ if and only if $\times a: M \to M$ is
an injection, or equivalently in the language of \S\emph{1.5}, if
$$
0 \longrightarrow M \xrightarrow{\times a} M
$$
is exact.)

\begin{proof}
	\begin{itemize}
        \item $(\Leftarrow)$: When $S$ contains no zerodivisors, we prove by
            contradiction: Assume that under the canonical ring map
            $\varphi:A\mapsto S^{-1}A$, there exist $a_{1},a_{2}\in A,\
            s.t.\varphi(a_{1})=\varphi(a_{2})$. Now we have $a_{1}/1 = a_{2}/1$,
            which means
			$$
			\exists s \in S, s.t.\ s(a_{1}-a_{2})=0
			.$$ 
			Which contradicts with the fact that $S$ has no zerodivisors!
        \item $(\Rightarrow)$: When $\varphi:A \mapsto  S^{-1}A$ is injective,
            we assume $S$ contains a zerodivisors $s$ at which
			$$
			\exists d\in A,\ s.t.\ d\neq 0,\ s\cdot d=0
			.$$ 
            Now let $a_{1},a_{2}\in A,\ s.t.\ a_{1}-a_{2}=d$. Since $\varphi$ is
            injective,  we must have $\varphi(a_{1})\neq \varphi(a_{2})$. Which
            contradicts with the fact that
			$$
			s(a_{1}-a_{2})=s\cdot d = 0
			.$$ 
	\end{itemize}
\end{proof}

\end{exercise}

\begin{exercise}
Verify that $A \rightarrow S^{-1}A$ satisfies the following universal
property: $S^{-1}A$ is initial among A-algebras $B$ where every element of
$S$ is sent to an invertible element in $B$. (Recall: the data of ``an
A-algebra $B$'' and ``a ring map $A \rightarrow B$'' are the same.)
Translation: any map $A \rightarrow B$ where every element of $S$ is sent to
an invertible element must factor uniquely through $A \rightarrow S^{-1}A$.
Another translation: a ring map out of $S^{-1}A$ is the same thing as a ring
map from $A$ that sends every element of $S$ to an invertible element.
Furthermore, an $S^{-1}A$-module is the same thing as an $A$-module for
which $s \times \cdot : M \rightarrow M$ is an $A$-module isomorphism for
all $s \in S$.

\begin{proof}
	Suppose
	\begin{align*}
		& f:A\mapsto B\\
		& \varphi :A\mapsto S^{-1}A
	.\end{align*}
    By the first translation in by author, we have to prove that there an unique
    ring homomorphism $\psi$ mapping $S^{-1}A$ to $B$. Recall that $\forall a
    \in A, f(a)=\varphi(a / 1)$. Also, for all $a\in A,s \in S$,
	$$
	\frac{a}{s}\cdot \frac{s}{1}= \frac{a}{1}
	.$$ 
	So
	\begin{align*}
		& \psi( \frac{a}{s} )\cdot \psi( \frac{s}{1} ) = \psi( \frac{a}{1} )\\
        & \psi( \frac{a}{s} )= \psi( \frac{a}{1} )\cdot \psi( \frac{s}{1} )^{-1}
        \\
	.\end{align*}
	And we can see $\psi( \frac{a}{s} )$ is uniquely defined!
\end{proof}

\end{exercise}

\begin{exercise}
Show that $\phi \colon M \to S^{-1}M$ exists, by constructing something
satisfying the universal property. Hint: define elements of $S^{-1}M$ to be
of the form $m/s$ where $m \in M$ and $s \in S$, and $m_1/s_1 = m_2/s_2$ if
and only if for some $s \in S$, $s(s_2m_1 - s_1m_2) = 0$. Define the
additive structure by $(m_1/s_1) + (m_2/s_2) = (s_2m_1 + s_1m_2)/(s_1s_2)$,
and the $S^{-1}A$-module structure (and hence the $A$-module structure) is
given by $(a_1/s_1) \cdot (m_2/s_2) = (a_1m_2)/(s_1s_2)$.

\begin{proof}
	We follow the hint from author: Define
	$$
	S^{-1}M= \left\{ \frac{m}{s}\mid m \in M,s \in S \right\} 
	.$$ 
    And $S^{-1}M$ meet the requirement to become a module by: $\forall
    m_{1},m_{2}\in M,s_{1},s_{2}\in S:$
	$$
	\frac{m_{1}}{s_{1}} = \frac{m_{2}}{s_{2}} \iff \exists s \in S,\ s \left( s_{2}m_{1}-s_{1}m_{2} \right) =0\\
	.$$ 
	$$
    \frac{m_{1}}{s_{1}} + \frac{m_{2}}{s_{2}}
     = \frac{m_{1}s_{2}+m_{2}s_{1}}{s_{1}s_{2}}
	.$$ 
	And the $S^{-1}A$ module structure
	$$
    \forall a \in A,s_{1},s_{2}\in S,m\in M,\ \frac{a}{s_{1}}\cdot
    \frac{m}{s_{2}} = \frac{am}{s_{1}s_{2}}
	.$$ 
	Also, we define $M\Rightarrow S^{-1}M$ by
	$$
	m \longrightarrow \frac{m}{1}
	.$$ 
    Now for any $A$-module maps $M\rightarrow N$, we'll prove $\phi$ satisfies
    the universal proterty, i.e. the map $S^{-1}M\longrightarrow N$, we denote
    it by $\varphi$, is unique: $\forall \frac{m}{s}\in S^{-1}M$, notice:
	$$
	\begin{cases}
		\alpha(m) \in N \\
		\alpha(s)^{-1} \in N
	\end{cases}
	.$$ 
    Therefore we have $\varphi( \frac{m}{s} ) = \alpha(m) \alpha(s)^{-1}$. So
    the $\phi$ in the problem satisfy the uniqueness of universal property, so
    it exist.
\end{proof}

\end{exercise}

\begin{exercise}
\begin{enumerate}[label=({\alph*})]
\item Show that localization commutes with finite products, or
    equivalently, with finite direct sums. In other words, if
    $M_{1},\ldots,M_{n}$ are A-modules, describe an isomorphism (of
    A-modules, and of $S^{-1}A$-modules)
    $S^{-1}(M_{1}\times\cdots\times M_{n}) \to S^{-1}M_{1}\times\cdots\times
    S^{-1}M_{n}$.
\item Show that localization commutes with arbitrary direct sums.
\item Show that ``localization does not necessarily commute with infinite
    products'': the obvious map $S^{-1}(\prod_{i}M_{i}) \to \prod_{i}
    S^{-1}M_{i}$ induced by the universal property of localization is not
    always an isomorphism. (Hint:
    $(1,1/2,1/3,1/4,\ldots)\in\mathbb{Q}\times\mathbb{Q}\times\cdots$)
\end{enumerate}

\begin{proof}
	\begin{enumerate}
		\item Intuitively, we can let
			$$
            \frac{(m_{1},m_{2},\ldots ,m_{n})}{s} \longrightarrow \left(
            \frac{m_{1}}{s},\ \frac{m_{2}}{s},\ \ldots,\ \frac{m_{n}}{s} \right) 
			.$$ 
            On the other hand, inductively we only have to consider the case
            $n=2$, and we let
			$$
            \left( \frac{m_{1}}{s_{1}},\ \frac{m_{2}}{s_{2}} \right)
            \longrightarrow \left( \frac{m_{1}s_{2}}{s_{1}s_{2}},\
            \frac{m_{2}s_{1}}{s_{1}s_{2}} \right) 
			.$$ 
			And we're done after induction.
        \item Notice that in the case of direct sum, all but finitely many
            coordinates is zero. Therefore the proof is exactly the same as the
            case in \emph{(a)}.
        \item Counterexample: consider group of modules
            $\{M_{n}\}_{n=1}^{\infty}$, and set of prime numbers
            $P=\{p_{n}\}_{n=1}^{\infty}$. We consider fraction of
            $\{M_{n}\}_{n=1}^{\infty}$ by positive integers, we cannot map
            ${\frac{m_1}{p_1},\frac{m_2}{p_2},...}$ to a proper element in
            $S^{-1}\left( \bigotimes _{k \in I}M_{k} \right) $ since $p_{1},\
            p_{2},\ \ldots$ cannot have a shared denominator.
	\end{enumerate}
\end{proof}

\end{exercise}

\begin{exercise}
EXERCISE (IF YOU HAVEN'T SEEN TENSOR PRODUCTS BEFORE).

Show that $\mathbb{Z}/(10) \otimes_{\mathbb{Z}} \mathbb{Z}/(12) \cong
\mathbb{Z}/(2)$. (This exercise is intended to give some hands-on practice
with tensor products.)

\begin{proof}
	More generally, we'll prove that
	$$
    \forall m,\ n \in \mathbb{Z}^{+},\ \mathbb{Z} / (m)
    \otimes_{\mathbb{Z}}\mathbb{Z} \ (n) \cong \mathbb{Z} / (\gcd(m,n))
	.$$ 
    First, we know $\mathbb{Z} / (m) \otimes_{\mathbb{Z}}\mathbb{Z} \ (n)$ is
    generated by $1_{\mathbb{Z} / (m)} \otimes 1_{\mathbb{Z} / (n)}$. Which means
    that every elements of $\mathbb{Z} / (m) \otimes_{\mathbb{Z}}\mathbb{Z} \
    (n) $ can be written in the form of $k\cdot (1_{\mathbb{Z} / (m)} \otimes
    1_{\mathbb{Z} / (n)}),\ k \in \mathbb{Z}^{+}$. Also, $\forall i,\ j \in
    \mathbb{Z}^{+}$
	,$$
    k\cdot (1_{\mathbb{Z} / (m)} \otimes 1_{\mathbb{Z} / (n)}) = (mi+nj)\cdot
    k\cdot (1_{\mathbb{Z} / (m)} \otimes 1_{\mathbb{Z} / (n)})
	.$$ 
	And the proof can be done by Euclidean algorithm.
\end{proof}

\end{exercise}

\begin{exercise}[Proved in \emph{Atiyah}]
IMPORTANT EXERCISE: RIGHT-EXACTNESS OF $(\cdot)\otimes_{A}N$.

Show that $(\cdot)\otimes_{A}N$ gives a covariant functor
$\mod_{A}\to\mod_{A}$. Show that $(\cdot)\otimes_{A}N$ is a right-exact
functor, i.e., if
$$
M^{\prime} \to M \to M^{\prime\prime} \to 0
$$
is an exact sequence of $A$-modules (which means $f\colon M\to M''$ is
surjective, and $M'$ surjects onto the kernel of $f$; see \S\emph{1.5}),
then the induced sequence
$$
M^{\prime} \otimes_{A} N \to M \otimes_{A} N \to M^{\prime\prime}
\otimes_{A} N \to 0
$$
is also exact. This exercise is repeated in Exercise 1.5.G, but you may get
a lot out of doing it now. (You will be reminded of the definition of
right-exactness in \S\emph{1.5.6}.)
\end{exercise}

\begin{exercise}[Trivial]
Show that $(T, t: M \times N \rightarrow T)$ is unique up to unique
isomorphism. Hint: first figure out what ``unique up to unique isomorphism''
means for such pairs, using a category of pairs $(T, t)$. Then follow the
analogous argument for the product.

\begin{proof}
    Because from the textbook content right above this question, we see any two
    maps $T,\ T'$ has a unique $A$-linear map to each other.
\end{proof}

\end{exercise}

\begin{exercise}
Show that the construction of \S\emph{1.2.5} satisfies the universal
property of tensor product.

\begin{proof}
    Because the construction in \S\emph{1.2.5} is a bilinear map, is unique up
    to isomorphism.
\end{proof}

\end{exercise}

\begin{exercise}
IMPORTANT EXERCISE.

\begin{enumerate}[label=({\alph*})]
\item If $M$ is an $A$-module and $A \to B$ is a morphism of rings, give
    $B \otimes_A M$ the structure of a $B$-module (this is part of the
    exercise). Show that this describes a functor $\mod_A \to \mod_B$.
\item (tensor product of rings) If further $A \to C$ is another morphism of
    rings, show that $B \otimes_A C$ has a natural structure of a ring.
    Hint: multiplication will be given by
    $(b_1 \otimes c_1)(b_2 \otimes c_2) = (b_1 b_2) \otimes (c_1 c_2)$.
    (Exercise 1.2.U will interpret this construction as a fibered
    coproduct.)
\end{enumerate}

\begin{proof}
	\begin{enumerate}
		\item 
			\begin{itemize}
				\item $\forall b_{1},\ b_{2}\in B,\ b\otimes_{A}m\in B\otimes M$, we have
					\begin{align*}
                        (b_{1}+b_{2})(b\otimes_{A}m)
                        &=(b_{1}+b_{2})b\otimes_{A}m\\
                        &=(b_{1}b+b_{2}b)\otimes_{A}m\\
                        &=b_{1}b\otimes_{A}m+b_{2}b\otimes_{A}m\\
                        &=b(b_{1}\otimes_{A}m)+b(b_{2}\otimes_{A}m)
					.\end{align*}
					\begin{align*}
                        (b_{1}b_{2})(b\otimes_{A}m)
                        &=(b_{1}b_{2}b)\otimes_{A}m\\
                        &=b_{1}(b_{2}(b\otimes_{A}m))
					.\end{align*}
                    $$
                        b(b_{1}\otimes_{A}m+b_{2}\otimes_{A}m)
                        =b(b_{1}\otimes_{A}m)+b(b_{2}\otimes_{A}m)
					.$$
                    % Which means that $B\otimes_{A}M$ meets the requirement to
                    % be a $B$-module.
				\item Denote such map by $\mathcal{F}$. We see for ring $A$ and
                    $B\otimes_{A}M\in \mod_{B}$,
					$$
					 \mathcal{F}: A\longmapsto B\otimes_{A}M
					.$$ 
					And for $\varphi \in \mor(A),\ \varphi: A\mapsto A'$,
					$$
					\mathcal{F}(\varphi):B\otimes_{A}M\longmapsto B\otimes_{A'}M
					.$$ 
			\end{itemize}
		\item We define addtion as formal sum and multiplication as
			$$
			(b_1\otimes c_1)\cdot (b_2\otimes c_2)=(b_1b_2)\otimes (c_1c_2)
			.$$ 
			And we can easily exam the properties to let it becomes a ring.
	\end{enumerate}
\end{proof}

\end{exercise}

\begin{exercise}
IMPORTANT EXERCISE.

If $S$ is a multiplicative subset of $A$ and $M$ is an $A$-module, describe
a natural isomorphism $(S^{-1}A) \otimes_A M \stackrel{\sim}{\longrightarrow}
S^{-1}M$ (as $S^{-1}A$-modules and as $A$-modules).

\begin{proof}
	We can define
	$$
	(S^{-1}A)\otimes _{A}M\longrightarrow S^{-1}M
	.$$ 
	As
	$$
	(\frac{a}{s})\otimes m\longmapsto \frac{am}{s}
	.$$ 
    And we can easily check that it meets the requirement to become a
    isomorphism.
\end{proof}

\end{exercise}

\begin{exercise}
EXERCISE ($\otimes$ COMMUTES WITH $\oplus$).

Show that tensor products commute with arbitrary direct sums: if $M$ and
$\{N_{i}\}_{i\in I}$ are all A-modules, describe an isomorphism
$$
M \otimes (\oplus_{i \in I} N_{i}) \xrightarrow{\sim} \oplus_{i \in I}
(M \otimes N_{i}).
$$

\begin{proof}
    Because $I$ only has finite nonzero elements, so the properties of tensor
    product can be verified easily.
\end{proof}

\end{exercise}

\begin{exercise}
EXERCISE (FIBERED PRODUCTS OF SETS).

Show that in Sets,
$$
X \times_{Z} Y = \{(x, y) \in X \times Y: \alpha(x) = \beta(y)\}.
$$
More precisely, show that the right side, equipped with its evident maps to
$X$ and $Y$, satisfies the universal property of the fibered product. (This
will help you build intuition for fibered products.)

\begin{proof}
	It's because
	$$
	\alpha\circ\pr_{X} = \beta\circ\pr_{Y}
	.$$ 
\end{proof}

\end{exercise}

\begin{exercise}
If $X$ is a topological space, show that fibered products always exist in
the category of open sets of $X$, by describing what a fibered product is.
(Hint: it has a one-word description.)

\begin{solution}
    Intersection.

    Because morphisms are inclusion embeddings.
\end{solution}

\end{exercise}

\begin{exercise}
If $Z$ is the final object in a category $\mathcal{C}$, and
$X, Y \in \mathcal{C}$, show that ``$X \times_{Z} Y = X \times Y$'': ``the''
fibered product over $Z$ is uniquely isomorphic to ``the'' product. Assume
all relevant (fibered) products exist. (This is an exercise about unwinding
the definition.)

\begin{proof}
	Denote the natural maps
	$$
	\pr_{X}:X\times Y\longrightarrow X
	.$$ 
	$$
	pr_{Y}:X\times Y\longrightarrow Y
	.$$ 
	And the unique map
	$$
	\alpha:X\longrightarrow Z
	.$$ 
	$$
	\beta:Y\longrightarrow Z
	.$$ 
	Now we have to prove
	$$
	\alpha\circ\pr_{X}=\beta\circ\pr_{Y}
	.$$ 
    Which do work since $Z$ is a final diagram, from which we know, up to
    isomorphisms, has only a single morphism sending $X\times Y$ to $Z$.
\end{proof}

\end{exercise}

\begin{exercise}
USEFUL EXERCISE: TOWERS OF CARTESIAN DIAGRAMS ARE CARTESIAN DIAGRAMS.

If the two squares in the following commutative diagram are Cartesian
diagrams, show that the ``outside rectangle'' (involving $U$, $V$, $Y$, and
$Z$) is also a Cartesian diagram.
$$
\begin{tikzcd}
    U \arrow[r] \arrow[d] & V \arrow[d] \\
    W \arrow[r] \arrow[d] & X \arrow[d] \\
    Y \arrow[r] & Z
\end{tikzcd}
$$

\begin{proof}
	For notations, let
	$$
	\begin{tikzcd}
        U \arrow[r, "f_1"] \arrow[d, "g_1"] & V \arrow[d, "h_1"] \\
        W \arrow[d, "g_2"] \arrow[r, "f_2"] & X \arrow[d, "h_2"] \\
        Y \arrow[r, "f_3"]                  & Z                 
    \end{tikzcd}
	.$$ 
	And we can see
	\begin{align*}
		h_2\circ h_1\circ f_1
        &=h_2\circ f_2\circ g_1\\
		&=f_3\circ g_2\circ g_1
	.\end{align*}
\end{proof}

\end{exercise}

\begin{exercise}
Given morphisms $X_{1} \rightarrow Y$, $X_{2} \rightarrow Y$, and
$Y \rightarrow Z$, show that there is a natural morphism $X_{1} \times_{Y}
X_{2} \rightarrow X_{1} \times_{Z} X_{2}$, assuming that both fibered
products exist. (This is trivial once you figure out what it is saying. The
point of this exercise is to see why it is trivial.)

\begin{proof}
	$$
	\begin{tikzcd}
        X_1\times_Y X_2 \arrow[rdd] \arrow[rrd] \arrow[rd, "\exists!" description, dotted] &                                       &                                   &   \\
                                                                                           & X_1\times_{Z} X_2 \arrow[d] \arrow[r] & X_2 \arrow[d] \arrow[rdd, dashed] &   \\
                                                                                           & X_1 \arrow[r] \arrow[rrd, dashed]     & Y \arrow[rd]                      &   \\
                                                                                           &                                       &                                   & Z
	\end{tikzcd}
	.$$ 
	Because from the problem we have the commutative diagram
	$$
	\begin{tikzcd}
        X_1\times_Y X_2 \arrow[d] \arrow[r] & X_2 \arrow[d, dashed] \\
        X_1 \arrow[r, dashed]               & Z                    
	\end{tikzcd}
	.$$ 
    And by the definition of fibered product, there must exist a unqiue
    factoring map $X_1\times_{Y}X_2\longrightarrow X_{1}\times _{Z}X_2$.
\end{proof}

\end{exercise}

\begin{exercise}
IMPORTANT EXERCISE: THE DIAGONAL-BASE-CHANGE DIAGRAM.

Suppose we are given morphisms $X_{1}, X_{2} \rightarrow Y$ and
$Y \rightarrow Z$. Show that the following diagram is a Cartesian square.
$$
\begin{tikzcd}
    X_1\times_Y X_2 \arrow[r] \arrow[d] & X_1\times_Z X_2 \arrow[d] \\
    Y \arrow[r] & Y\times_Z Y
\end{tikzcd}
$$
Assume all relevant (fibered) products exist. (If this exercise is too hard
now, you can try it again at Exercise 1.3.B.)

\begin{proof}
    \emph{(}$ Y\longrightarrow Y\times _{Z}Y $\emph{is defined by identical
    embedding)}
	$$
	\begin{tikzcd}
X_1\times_Y X_2 \arrow[d] \arrow[r] & X_2 \arrow[d] &   \\
X_1 \arrow[r]                       & Y \arrow[rd]  &   \\
                                    &               & Z
	\end{tikzcd}
	$$ 
	We need to prove the diagram below is a Cartesian Square:
	$$
	\begin{tikzcd}
        X_1\times_Y X_2 \arrow[d] \arrow[r] & X_1\times_Z X_2 \arrow[d] \\
        Y \arrow[r]                         & Y\times_Z Y              
	\end{tikzcd}
	$$ 
    Which becomes trivial when we notice
    $$
    \begin{tikzcd}
        X_1\times_{Y}X_2 \arrow[rd, "\exists!" description, dotted] \arrow[rrd] \arrow[rdd] &                                                                                 &                                             &             \\
                                                                                        & X_1\times_{Z}X_2 \arrow[rd, "\exists!" description, dotted] \arrow[d] \arrow[r] & X_2 \arrow[rd]                              &             \\
                                                                                        & X_1 \arrow[rd]                                                                  & Y\times_{Z}Y \arrow[d] \arrow[r]            & Y \arrow[d] \\
                                                                                        &                                                                                 & Y \arrow[ru, Rightarrow,no head] \arrow[r] & Z          
    \end{tikzcd}
    $$
\end{proof}

\end{exercise}

\begin{exercise}
Show that coproduct for Sets is disjoint union. This is why we use the
notation $\coprod$ for disjoint union.

\begin{proof}
	For
	$$
	\begin{tikzcd}
                                                & S                                                       &                                                  \\
S_1 \arrow[r, "\varphi_1"] \arrow[ru, "\psi_1"] & S_1\sqcup S_2 \arrow[u, dotted] & S_2 \arrow[l, "\varphi_2"] \arrow[lu, "\psi_2"']
	\end{tikzcd}
	.$$ 
	For $s_1\in S_1,\ s_2\in S_2$, define
	$$
	\varphi_1(s_1)=(s_1,\cdot ),\ \varphi_2(s_2)=(\cdot ,s_2)
	.$$
	Now for any $\psi_1,\ \psi_2$, we naturally have the induced morphism
	$$
	\Psi :S_1\sqcup S_2\longrightarrow S
	$$ 
	$$
	\Psi:(s_1,\cdot )\longmapsto \psi_1(s_1)
	$$ 
	$$
	\Psi:(\cdot ,s_2)\longmapsto \psi_2(s_2)
	.$$
	Which is unique up to isomorphism while making the diagram commutes.
\end{proof}

\end{exercise}

\begin{exercise}
Suppose $A \to B$ and $A \to C$ are two ring morphisms, so in particular $B$
and $C$ are $A$-modules. Recall (Exercise 1.2.K) that $B \otimes_A C$ has a
ring structure. Show that there is a natural morphism $B \to B \otimes_A C$
given by $b \mapsto b \otimes 1$. (This is not necessarily an inclusion; see
Exercise 1.2.G.) Similarly, there is a natural morphism $C \to B \otimes_A
C$. Show that this gives a fibered coproduct on rings, i.e., that
$$
\begin{tikzcd}
    A \arrow[r] \arrow[d] & B \arrow[d] \\
    C \arrow[r] & B \otimes_A C
\end{tikzcd}
$$
satisfies the universal property of fibered coproduct.

\begin{proof}
	Denote
	$$
	\begin{tikzcd}
D &                                                       &                               \\
  & B\otimes_AC \arrow[lu, "\varphi" description, dotted] & C \arrow[l] \arrow[llu, "h"'] \\
  & B \arrow[u] \arrow[luu, "g"]                          & A \arrow[l] \arrow[u]        
	\end{tikzcd}
	$$ 
	and we obtain the unique map $\varphi$ by
	$$
	\varphi((a\cdot b)\otimes c)=g(b)
	$$ 
	$$
	\varphi(b\otimes (a\cdot c))=h(c)
	$$ 
\end{proof}

\end{exercise}

\begin{exercise}
Show that the composition of two monomorphisms is a monomorphism.

\begin{proof}
	$$
	\begin{tikzcd}
        A \arrow[d, "\leq1" description, dotted] \arrow[rd, "\leq1" description, dotted] \arrow[rrd] &                       &   \\
        X \arrow[r, "\alpha"']                                                                       & Y \arrow[r, "\beta"'] & Z
	\end{tikzcd}
	$$ 
	From the graph we obtain the proof easily.
\end{proof}

\end{exercise}

\begin{exercise}
Prove that a morphism $\pi\colon X\to Y$ is a monomorphism if and only if
the fibered product $X\times_{Y}X$ exists, and the induced diagonal morphism
$\delta_{\pi}\colon X\to X\times_{Y}X$ (Definition 1.2.7) is an
isomorphism. We may then take this as the definition of monomorphism.
(Monomorphisms aren't central to future discussions, although they will come
up again. This exercise is just good practice.)

\begin{proof}
    \begin{itemize}
    	\item 
            When $\pi$ is a monomorphism:
            For
            $$
            \begin{tikzcd}
                X\times_{Y}X \arrow[r] \arrow[d] & X \arrow[d, "\pi"] \\
                X \arrow[r, "\pi"]               & Y                 
            \end{tikzcd}
            $$
            From the definition of monomorphism, we know there's only one
            morphism in $\mor{(X\times _{Y}X,X)}$ making the diagram commutes.
            Therefore the fiber product must exist, where the morphism in
            between is an isomorphism.
        \item When $X\times _{Y}X$ exists, then any $A$ letting the diagram
            below commute
    		$$
    		\begin{tikzcd}
    A \arrow[rdd] \arrow[rrd] \arrow[rd, "\exists!" description, dotted] &                                &                    \\
                                                         & X\times_YX \arrow[d] \arrow[r] & X \arrow[d, "\pi"] \\
                                                         & X \arrow[r, "\pi"']            & Y                 
    	    \end{tikzcd}
    	,$$ 
        whose morphism to $X\times _{Y}X$ will factor through $X\times _{Y}X$
        uniquely. Therefore any two map:
    	$$
    	\alpha,\ \beta:A\longrightarrow X
    	$$ 
    	can be factorized into
    	$$
    	\alpha= \alpha'\circ p,\ \beta=\beta'\circ p
    	$$ 
        Where $p:A\longrightarrow X\times _{Y}X$ and $\alpha',\
        \beta':X\times_{Y}X\longrightarrow X$. \\
    	$$
    \begin{tikzcd}
            A \arrow[rdd, "\alpha"'] \arrow[rrd, "\beta"] \arrow[rd, "p" description, dotted] &                                                      &                    \\
                                                                                              & X\times_YX \arrow[d, "\alpha'"'] \arrow[r, "\beta'"] & X \arrow[d, "\pi"] \\
                                                                                              & X \arrow[r, "\pi"']                                  & Y                 
    \end{tikzcd}
    $$ 
    So if
    $$
    \pi\circ \alpha=\pi\circ \beta
    $$ 
    We'll have
    $$
    \pi\circ \alpha'\circ p=\pi\circ \beta'\circ p
    $$ 
    By the definiton of diagonal morphism we know $\alpha'=\beta'$, therefore
    $\alpha'\circ p=\beta'\circ p$, i.e. $\alpha=\beta$. $\pi$ satisfy the left
    cancellation law, so it's a monomorphism.
	\end{itemize}
\end{proof}

\end{exercise}

\begin{exercise}
EASY EXERCISE.

We use the notation of Exercise 1.2.R. Show that if $Y \rightarrow Z$ is a
monomorphism, then the morphism $X_{1} \times_{Y} X_{2} \rightarrow X_{1}
\times_{Z} X_{2}$ you described in Exercise 1.2.R is an isomorphism. (Hint:
for any object $V$, give a natural bijection between maps from $V$ to the
first and maps from $V$ to the second. It is also possible to use the
Diagonal-Base-Change diagram, Exercise 1.2.S.)

\begin{proof}
	$$
	\begin{tikzcd}
        X_1\times_YX_2 \arrow[rdd] \arrow[rrd] &                                    &               &   \\
                                               & X_1\times_ZX_2 \arrow[d] \arrow[r] & X_2 \arrow[d] &   \\
                                               & X_1 \arrow[r]                      & Y \arrow[rd]  &   \\
                                               &                                    &               & Z
	\end{tikzcd}
	$$ 
	From Exercise \textbf{1.2.R} we know there exist a unique morphism
	$$
	\varphi:X_1\times_{Y}X_2\longrightarrow X_1\times _{Z}X_2
	.$$ 
    Now since $Y\longrightarrow Z$ is monomorphism, from the commutavity of
    $X_1\times _{Z}X_2\rightarrow X_1\rightarrow Y\rightarrow Z$ and $X_1\times
    _{Z}X_2\rightarrow X_2\rightarrow Y\rightarrow Z$ we can get the commutavity
    of
	$$
	\begin{tikzcd}
X_1\times_ZX_2 \arrow[d] \arrow[r] & X_2 \arrow[d] \\
X_1 \arrow[r]                      & Y            
	\end{tikzcd}
	$$ 
    And therefore there's a unique morphism $\psi:X_1\times
    _{Z}X_2\longrightarrow X_1\times _{Y}X_2$. By their uniqueness, we have
	$$
	\varphi\circ \psi\circ \varphi=\varphi,\ \psi\circ \varphi\circ \psi=\psi
	$$ 
    So $\varphi\circ \psi=\id_{X_1\times _{Z}X_2},\ \psi\circ
    \varphi=\id_{X_1\times _{Y}X_2}$.
\end{proof}

\end{exercise}

\begin{exercise}
IMPORTANT EXERCISE THAT YOU SHOULD DO ONCE IN YOUR LIFE (YONEDA'S LEMMA).

\begin{enumerate}[label=({\alph*})]
\item Suppose you have two objects $A$ and $A'$ in a category $\mathcal{C}$,
    and maps
    $$
    i_{C} \colon \mor(C, A) \longrightarrow \mor(C, A')
    $$
    that commute with the maps (1.2.11.1). Show that the $i_C$ (as $C$
    ranges over the objects of $\mathcal{C}$) are induced from a unique
    morphism $g\colon A\to A'$. More precisely, show that there is a unique
    morphism $g \colon A \to A'$ such that for all $C \in \mathcal{C}$,
    $i_C$ is $u \mapsto g \circ u$.
\item If furthermore the $i_C$ are all bijections, show that the resulting
    $g$ is an isomorphism. (Hint for both: This is much easier than it
    looks. This statement is so general that there are really only a couple
    of things that you could possibly try. For example, if you're hoping to
    find a morphism $A \to A'$, where will you find it? Well, you are
    looking for an element $\mor(A, A')$. So just plug in $C = A$ to
    (1.2.11.2), and see where the identity goes.)
\end{enumerate}

\begin{proof}
	$$
	\begin{tikzcd}
B \arrow[rrdd] \arrow[dd]                                                    &                                                                               & A' \arrow[lldd, "h"', bend right, shift right=3] \arrow[loop, distance=2em, in=125, out=55] \\
{}                                                                           &                                                                               & {}                                                                                          \\
A \arrow[rruu, "g"', bend right] \arrow[loop, distance=2em, in=305, out=235] & {} \arrow[ru, "i_c"', bend right=67, shift right=2] \arrow[lu, bend right=49] & C \arrow[uu] \arrow[ll, "u" description]                                                   
	\end{tikzcd}
	$$ 
	\begin{enumerate}
        \item We simply let $\id_{C}$ correspond to the morphism
            $A\longrightarrow A'$ and verify the properties stated in the
            problems. Also, uniqueness can be obtained by letting $C=A$.
		\item We follow from the hint: Let $C=A'$ and we get a bijection
			$$
			\id_{A'}:(\mor(A',A))\longrightarrow (\mor(A',A'))
			$$ 
            We pick out $h$ from $\mor{(A',A)}$ such that
            $\id_{A'}(h)=\id_{A'}$. From \textbf{(a)} we know $\id_{A'}$ can be
            represented using $g$ :
			$$
			\id_{A'}=\id_{A'}(h)=g\circ h
			$$ 
            On the other hand, we prove $\id_{A}^{-1}$ can be induced by $h$
            using the similar method to \textbf{(a)}.
			$$
            \exists !\ \id_{A}\in \mor{(A,A)},\ \id_{A}(\id_{A})=g\in \mor{(A,A)}
			$$ 
			And we can see
			$$
			\id_{A}=\id^{-1}(g)=h\circ g
			$$ 
	\end{enumerate}
\end{proof}

\end{exercise}

\begin{exercise}[Skipped temporarily due to the star.]
$\ast$EXERCISE.

\begin{enumerate}[label=({\alph*})]
\item Suppose $A$ and $B$ are objects in a category $\mathcal{C}$. Give a
    bijection between the natural transformations $h^A \to h^B$ of covariant
    functors $\mathcal{C} \to Sets$ (see Example 1.1.14 for the definition)
    and the morphisms $B \to A$.
\item State and prove the corresponding fact for contravariant functors
    $h_{A}$ (see Example 1.1.20). Remark: A contravariant functor $F$ from
    $\mathcal{C}$ to Sets is said to be representable if there is a natural
    isomorphism
    $$
    \xi \colon F \xrightarrow{\sim} h_{A}.
    $$
    Thus the representing object $A$ is determined up to unique isomorphism
    by the pair $(F, \xi)$. There is a similar definition for covariant
    functors. (We will revisit this in \S\emph{7.6}, and this problem will
    appear again as Exercise 7.6.C. The element $\xi^{-1}(\mathrm{id}_A) \in
    F(A)$ is often called the ``universal object''; do you see why?)
\item Yoneda's Lemma. Suppose $F$ is a covariant functor $\mathcal{C} \to
    Sets$, and $A \in \mathcal{C}$. Give a bijection between the natural
    transformations $h^A \to F$ and $F(A)$. (The corresponding fact for
    contravariant functors is essentially Exercise 10.1.B.)
\end{enumerate}

\begin{proof}
	\begin{enumerate}
        \item We consider the contravariant functors of the morphisms
            $B\longrightarrow A$.
		\item 
	\end{enumerate}
\end{proof}

\end{exercise}

\section{Limits and colimits}
\begin{exercise}[Trivial]
EXERCISE (REALITY CHECK).

Suppose that the partially ordered set $I$ has an initial object $e$. Show
that the limit of any diagram indexed by $I$ exists.

\begin{proof}
    Since $\mathscr{I}$ is a partially ordered set, there's at most one morphism
    for any two object in $\mathscr{I}$. 
\end{proof}

\end{exercise}

\begin{exercise}
EXERCISE: THE DIAGONAL-BASE-CHANGE DIAGRAM, AGAIN.

Solve 1.2.S again by identifying both $X_{1} \times_{Y} X_{2}$ and
$Y \times_{(Y \times_{Z} Y)} (X_{1} \times_{Z} X_{2})$ as the limit of the
diagram
$$
\begin{tikzcd}
    X_1 \arrow[rd] & & X_2 \arrow[ld] \\
    & Y \arrow[d] & \\
    & Z &
\end{tikzcd}
$$

\begin{proof}
	$$
	\begin{tikzcd}
        X_1\times_ZX_2 \arrow[rd, dotted] \arrow[rdd] \arrow[rrd] &                                    &                                   &   \\
                                                                  & X_1\times_YX_2 \arrow[r] \arrow[d] & X_2 \arrow[d] \arrow[rdd, dashed] &   \\
                                                                  & X_1 \arrow[r] \arrow[rrd, dashed]  & Y \arrow[rd]                      &   \\
                                                                  &                                    &                                   & Z
	\end{tikzcd}
	\begin{tikzcd}
        Y\times_{(Y\times_ZY)}(X_1\times_ZX_2) \arrow[d] \arrow[r] & X_1\times_ZX_2 \arrow[d] \\
        Y \arrow[r]                                                 & Y\times_ZY              
	\end{tikzcd}
	$$ 
	For $X_1\times _{Y}X_2$:
	\begin{align*}
		X_1\times _{Y}X_2\rightarrow Y
        &=X_1\times _{Y}X_2\rightarrow X_1\rightarrow Y\\
		&=X_1\times _{Y}X_2\rightarrow X_2\rightarrow Y
	\end{align*}
	$$
	X_1\times _{Y}X_2\rightarrow Z=X_1\times _{Y}X_2\rightarrow Y\rightarrow Z
	$$ 
	For $Y\times_{(Y\times _{Z}Y)}(X_1\times _{Z}X_2)$:
	\begin{align*}
		Y\times_{(Y\times _{Z}Y)}(X_1\times _{Z}X_2)\rightarrow Y
		&=Y\times_{(Y\times _{Z}Y)}(X_1\times _{Z}X_2)\rightarrow X_1\times _{Z}X_2\rightarrow X_1\rightarrow Y\\
		&=Y\times_{(Y\times _{Z}Y)}(X_1\times _{Z}X_2)\rightarrow X_1\rightarrow Y
	\end{align*}
	$$
	Y\times _{(Y\times _{Z}Y)}(X_1\times _{Z}X_2)\rightarrow Z=
	Y\times _{(Y\times _{Z}Y)}(X_1\times _{Z}X_2)\rightarrow Y\rightarrow Z
	$$ 
\end{proof}

\end{exercise}

\begin{exercise}
IMPORTANT EXERCISE.

Show that in the category Sets,
$$
\left\{(a_{i})_{i \in \mathscr{I}} \in \prod_{i} A_{i}: F(m)(a_{j}) =
a_{k} \text{ for all } m \in \mor_{\mathscr{I}}(j, k) \subset
\mor(\mathscr{I})\right\},
$$
along with the obvious projection maps to each $A_{i}$, is the limit
$\lim_{\mathcal{I}} A_{i}$.

This clearly also works in the category $\mod_{A}$ of A-modules (in
particular $Vec_{k}$ and Ab), as well as Rings.

\begin{proof}
	$$
	\begin{tikzcd}
        (a_i)_{i\in \mathscr{I}} \arrow[d] \arrow[rd] &     \\
        A_j \arrow[r]                       & A_k
    \end{tikzcd}
	$$ 
	Because from the definition of $(a_{i})_{i\in \mathscr{I}}$, we have
	$$
	F(m)(a_{j})=a_{k}
	$$ 
    So $(a_{i})_{i\in \mathscr{I}}A\longrightarrow A_{j}\longrightarrow A_{k}$
    commutes with $(a_{i})_{i\in \mathscr{I}}\longrightarrow A_{k}$.
\end{proof}

\end{exercise}

\begin{exercise}
\begin{enumerate}[label=({\alph*})]
\item Interpret the statement ``$\mathbb{Q} = \text{colim} \frac{1}{n}
    \mathbb{Z}$''.
\item Interpret the union of some subsets of a given set as a colimit.
    (Dually, the intersection can be interpreted as a limit.) The objects of
    the category in question are the subsets of the given set.
\end{enumerate}

\begin{solution}
	\begin{enumerate}
        \item It means fractions converting into different denominators stands
            for the same rational number, which is obviously true.
        \item Because we can view the set category as posets, where morphisms in
            between two objects are unique.
	\end{enumerate}
\end{solution}

\end{exercise}

\begin{exercise}
Suppose $\mathcal{I}$ is filtered. (We will almost exclusively use the case
where $\mathcal{I}$ is a filtered set.) Recall the symbol $\coprod$ for
disjoint union of sets. Show that any diagram in Sets indexed by
$\mathcal{I}$ has the following, with the obvious maps to it, as a colimit:
$$
\left\{(a_{i}, i) \in \coprod_{i \in \mathscr{I}} A_{i}\right\} /
\{(a_{i}, i) \sim (a_{j}, j) \text{ if and only if there are } f \colon
A_{i} \to A_{k} \text{ and } g \colon A_{j} \to A_{k} \text{ in the diagram
for which } f(a_{i}) = g(a_{j}) \text{ in } A_{k}\}
$$
(You will see that the ``filtered'' hypothesis is there to ensure that
$\sim$ is an equivalence relation.)

\begin{proof}
	The proof is done by definition.
\end{proof}

\end{exercise}

\begin{exercise}
Verify that the A-module described above is indeed the colimit. (Make sure
you verify that addition is well-defined, i.e., is independent of the choice
of representatives $m_{i}$ and $m_{j}$, the choice of $\ell$, and the choice
of arrows $u$ and $v$. Similarly, make sure that scalar multiplication is
well-defined.)

\begin{proof}
    \emph{(Trivial, skipped.)}
\end{proof}

\end{exercise}

\begin{exercise}
USEFUL EXERCISE (LOCALIZATION AS A COLIMIT).

Generalize Exercise 1.3.D(a) to interpret localization of an integral
domain as a colimit over a filtered set: suppose $S$ is a multiplicative set
of $A$, and interpret $S^{-1}A = \text{colim} \frac{1}{s} A$ where the
colimit is over $s \in S$, and in the category of $A$-modules. (Aside: Can
you make some version of this work even if $A$ isn't an integral domain,
e.g., $S^{-1}A = \text{colim} A_s$? This will work in the category of
$A$-algebras.)

\begin{proof}
	\emph{(The proof is similar to \textbf{1.3.D} I'll skip.)}
\end{proof}

\end{exercise}

\begin{exercise}
EXERCISE: COLIMITS OF A-MODULES WITHOUT THE FILTERED CONDITION.

Suppose you are given a diagram of A-modules indexed by $\mathscr{I}$:
$F: \mathscr{I} \to Mod_A$, where we let $M_i := F(i)$. Show that the
colimit is $\oplus_{i \in \mathscr{I}} M_i$ modulo the relations
$m_i - F(n)(m_i)$ for every $n: i \to j$ in $\mathscr{I}$ (i.e., for every
arrow in the diagram). (Somewhat more precisely: ``modulo'' means
``quotiented by the submodule generated by''.)

\begin{proof}
    \emph{(Trivial, skipped.)}
    $$
        \begin{tikzcd}
        \colim_{i\in\mathscr{I}}{M_i} &                          \\
        M_i \arrow[u]                 & M_j \arrow[lu] \arrow[l]
        \end{tikzcd}
    $$
\end{proof}

\end{exercise}

\section{Adjoints}
\begin{exercise}
Write down what this diagram should be.

\begin{solution}
    $$
        \begin{tikzcd}
            {\mor_{\mathscr{B}}(F(A),B)} \arrow[r, "g^*"] \arrow[d, "\tau_{\mathcal{AB}}"] & {\mor_{\mathscr{B}}(F(A),B')} \arrow[d, "\tau_{\mathcal{AB'}}"] \\
            {\mor_{\mathscr{A}}(A,G(B))} \arrow[r, "Gg^*"]                                 & {\mor_{\mathscr{A}}(A,G(B'))}                                  
        \end{tikzcd}
    $$
\end{solution}

\end{exercise}

\begin{exercise}
Show that the map $\tau_{AB}$ (1.4.0.1) has the following properties. For
each $A$ there is a map $\eta_{A}: A \to \mathrm{GF}(A)$ so that for any
$g: F(A) \to B$, the corresponding $\tau_{AB}(g): A \to G(B)$ is given by
the composition
$$
A \xrightarrow{\eta_{A}} GF(A) \xrightarrow{Gg} G(B).
$$
Similarly, there is a map $\epsilon_{B}\colon FG(B)\to B$ for each $B$ so
that for any $f\colon A\to G(B)$, the corresponding map
$\tau_{AB}^{-1}(f)\colon F(A)\to B$ is given by the composition
$$
F(A) \xrightarrow{Ff} FG(B) \xrightarrow{\epsilon_{B}} B.
$$

\begin{proof}
    Use the natural property: For all $g:B'\longrightarrow B$, we have
    $$
        \begin{tikzcd}
        {\mor_{\mathscr{B}}{(F(A),B')}} \arrow[r, "g^*"] \arrow[d, "\tau_{AB'}"] & {\mor_{\mathscr{B}}{(F(A),B)}} \arrow[d, "\tau_{AB}"] \\
        {\mor_{\mathscr{A}}{(A,G(B'))}} \arrow[r, "Gg^*"]                        & {\mor_{\mathscr{A}}{(A,G(B))}}                       
        \end{tikzcd}
    $$
    commutes. Therefore let $B'=F(A)$ and we'll have
    $$
        \begin{tikzcd}
        {\mor_{\mathscr{B}}{(F(A),F(A))}} \arrow[r, "g^*"] \arrow[d, "\tau_{AF(A)}"] & {\mor_{\mathscr{B}}{(F(A),B)}} \arrow[d, "\tau_{AB}"] \\
        {\mor_{\mathscr{A}}{(A,GF(A))}} \arrow[r, "Gg^*"]                            & {\mor_{\mathscr{A}}{(A,G(B))}}                       
        \end{tikzcd}
    $$
    commutes. Take $\id_{F(A)}\in \mor_{\mathscr{B}}(F(A),F(A))$ we would have
    \begin{align*}
        \tau_{AB}(g)
        &= \tau_{AB}\circ g\circ \id_{F(A)} \\
        &= G(g)\circ \tau_{AF(A)}\circ \id_{F(A)}
    .\end{align*}
    Let $\eta_{A}=\tau_{AF(A)}(\id_{F(A)})$ and we're done.

    On the other hand: vice versa.
\end{proof}

\end{exercise}

\begin{exercise}
Suppose $M$, $N$, and $P$ are A-modules (where $A$ is a ring). Describe a
bijection $\hom_{A}(M \otimes_{A} N, P) \leftrightarrow \hom_{A}(M,
\hom_{A}(N, P))$. (Hint: try to use the universal property of $\otimes$.)

\begin{proof}
    For $\forall \beta\in \bil{\left( M,N;P \right) }$, we may define
    $$
    f:M\longrightarrow \hom_{A}{\left( N,P \right)}
    $$
    $$
    m\longmapsto \beta\left( m,\left( - \right)  \right) 
    $$
    i.e. $f\left( m \right) \left( n \right) =\beta\left( m,n \right) $.

    For $\forall f\in \hom_{A}\left( M,\hom_{A}\left( N,P \right)  \right) $,
    let
    $$
    \beta:M\times N\longrightarrow P
    $$
    $$
    \left( m,n \right) \longmapsto f\left( m \right) \left( n \right)
    $$
    i.e. $\beta\left( m,n \right) =f\left( m \right) \left( n \right)$.

    This defines the bijection.
\end{proof}

\end{exercise}

\begin{exercise}
EXERCISE (TENSOR-HOM ADJUNCTION).

Show that $(\cdot)\otimes_{A}N$ and $\hom_{A}(N,\cdot)$ are adjoint
functors.

\begin{proof}
    Define $F,G$ to be the functor $(\cdot)\otimes _{A}N,\hom_{A}(N,\cdot)$
    respectively. We need to prove those $\tau_{XY}$ are natual, i.e. for
    $\forall \varphi:M\longrightarrow M',\psi:P\longrightarrow P'$
    $$
    \begin{tikzcd}
        {\hom_{A}{(M'\otimes_{A}N,P)}} \arrow[r, "F\varphi^{*}"] \arrow[d, "\tau_{M'P}"] & {\hom_{A}{(M\otimes_{A}N,P)}} \arrow[d, "\tau_{MP}"] \\
        {\hom_{A}{(M',\hom_{A}{(N,P)})}} \arrow[r, "\varphi^{*}"]                        & {\hom_{A}{(M,\hom_{A}{(N,P)})}}                    
    \end{tikzcd}
    $$
    and
    $$
    \begin{tikzcd}
        {\hom_{A}{(M\otimes_{A}N,P)}} \arrow[r, "\psi^{*}"] \arrow[d, "\tau_{MP}"] & {\hom_{A}{(M\otimes_{A}N,P')}} \arrow[d, "\tau_{MP'}"] \\
        {\hom_{A}{(M,\hom_{A}{(N,P)})}} \arrow[r, "G\psi^{*}"]                     & {\hom_{A}{(M,\hom_{A}{(N,P')})}}                     
    \end{tikzcd}
    $$
    commutes.
    This is obvious when we consider the proeprty of homomorphisms.
\end{proof}

\end{exercise}

\begin{exercise}
Suppose $B \rightarrow A$ is a morphism of rings. If $M$ is an $A$-module,
you can create a $B$-module $M_{B}$ by considering it as a $B$-module. This
gives a functor $\cdot_{B}: Mod_{A} \rightarrow Mod_{B}$. Show that this
functor is right-adjoint to $\cdot \otimes_{B} A$. In other words, describe
a bijection
$$
\hom_{A}(N \otimes_{B} A, M) \cong \hom_{B}(N, M_{B})
$$
functorial in both arguments. (This adjoint pair is very important.)

\begin{proof}
    \begin{enumerate}[label=({\Roman*})]
        \item 
            $\tau_{NM}$:
            For $\Phi:N\otimes _{B}A\longrightarrow M$, define
            $$
            \phi:N\longrightarrow M_{B}
            $$
            $$
            n\longmapsto \Phi\left( n\otimes _{B}1_{A} \right) 
            $$
            
            $\tau_{NM}^{-1}$:
            On the other hand, for $\phi:N\longrightarrow M_{B}$, define
            $$
            \Phi:N\otimes _{B}A\longrightarrow M
            $$
            by
            $$
            \begin{tikzcd}
            N\times A \arrow[rd, "\phi\times\id"] \arrow[dd, "\pi"] &                          &   \\
                                                                    & M_{B}\times A \arrow[rd] &   \\
            N\times_{B}A \arrow[rr, "\Phi", dashed]                 &                          & M
            \end{tikzcd}
            $$
            and we have
            $$
            \Phi:n\otimes _{B}a\longmapsto a\phi\left( n \right) 
            $$

            Obviously the maps defined above are inverses to each other.
        \item 
            We need to show such bijection (transformation) is natural, i.e.
            $\forall f\in \hom_{A}\left( M,M' \right)$,
            $$
            \begin{tikzcd}
            {\hom_{A}(N\otimes_{B}A,M)} \arrow[d, "\tau_{NM}"] \arrow[r,"f^{*}"] & {\hom_{A}(N\otimes_{B}A,M')} \arrow[d, "\tau_{NM'}"] \\
            {\hom_{B}(N,M_{B})} \arrow[r, "(f^{*})_{B}"]                          & {\hom_{B}(N,M'_{B})}                               
            \end{tikzcd}
            $$
            commute, while $\forall g\in \hom_{B}\left( N,N' \right) $,
            $$
            \begin{tikzcd}
            {\hom_{A}(N'\otimes_{B}A,M)} \arrow[d, "\tau_{N'M}"] \arrow[r,"g^{*}\otimes_{B}\id_{A}"] & {\hom_{A}(N\otimes_{B}A,M)} \arrow[d, "\tau_{NM}"] \\
            {\hom_{B}(N',M_{B})} \arrow[r, "g^{*}"]                                                   & {\hom_{B}(N,M_{B})}                              
            \end{tikzcd}
            $$
            commutes.

            Pick a homomorphism from the start, then track it through the
            diagram directly and we're done.
    \end{enumerate}
\end{proof}

\end{exercise}

\begin{exercise}[Trivial]
EXERCISE (AN ABELIAN GROUP IS GROUPIFIED BY ITSELF).

Show that if an abelian semigroup is already a group then the identity
morphism is the groupification. (More correct: the identity morphism is a
groupification.) Note that you don't need to construct groupification (or
even know that it exists in general) to solve this exercise.
\end{exercise}

\begin{exercise}[Follow the hint, trivial]
Construct the ``groupification functor'' $H$ from the category of nonempty
abelian semigroups to the category of abelian groups. (One possible
construction: given an abelian semigroup $S$, the elements of its
groupification $H(S)$ are ordered pairs $(a, b) \in S \times S$, which you
may think of as $a - b$, with the equivalence that $(a, b) \sim (c, d)$ if
$a + d + e = b + c + e$ for some $e \in S$. Describe addition in this
group, and show that it satisfies the properties of an abelian group.
Describe the abelian semigroup map $S \to H(S)$.) Let $F$ be the forgetful
functor from the category of abelian groups Ab to the category of abelian
semigroups. Show that $H$ is left-adjoint to $F$.

\begin{solution}
    \emph{(It's just like constructing a fraction field of a ring.)}
\end{solution}

\end{exercise}

\begin{exercise}
EXERCISE (CF. EXERCISE 1.4.E).

The purpose of this exercise is to give you more practice with ``left
adjoints of forgetful functors'', the means by which we get abelian groups
from abelian semigroups, and sheaves from presheaves. Suppose $A$ is a ring,
and $S$ is a multiplicative subset. Then $S^{-1}A$-modules are a full
subcategory (\S\emph{1.1.15}) of the category of A-modules (via the obvious
inclusion $\mod_{S^{-1}A} \hookrightarrow \mod_{A}$). Then $\mod_{A} \to
\mod_{S^{-1}A}$ can be interpreted as an adjoint to the forgetful functor
$\mod_{S^{-1}A} \to \mod_{A}$. State and prove the correct statements.

(Here is the larger story. Every $S^{-1}A$-module is an $A$-module, and this
is an injective map, so we have a forgetful functor $F: \mod_{S^{-1}A} \to
\mod_{A}$. In fact this is a fully faithful functor: it is injective on
objects, and the morphisms between any two $S^{-1}A$-modules as $A$-modules
are just the same when they are considered as $S^{-1}A$-modules. Then there
is a functor $G: \mod_{A} \to \mod_{S^{-1}A}$, which might reasonably be
called ``localization with respect to $S$'', which is left-adjoint to the
forgetful functor. Translation: If $M$ is an $A$-module, and $N$ is an
$S^{-1}A$-module, then $\mor(GM, N)$ (morphisms as $S^{-1}A$-modules, which
are the same as morphisms as $A$-modules) are in natural bijection with
$\mor(M, FN)$ (morphisms as $A$-modules).)

\begin{proof}
    From \textbf{Exercise 1.4.E}, for ring homomorphism
    $$
    A\longrightarrow S^{-1}A
    $$
    we have
    $$
    \hom_{S^{-1}A}\left( S^{-1}N,M \right) \cong \hom_{A}\left( N,M_{A} \right) 
    $$
    \emph{(I suppose the rest is trivial.)}
\end{proof}

\end{exercise}

\section{An introduction to abelian categories}

\begin{exercise}[Trivial]
Describe exact sequences
$$
0 \longrightarrow \im f^{i} \longrightarrow A^{i+1} \longrightarrow
\coker f^{i} \longrightarrow 0
$$
$$
0 \longrightarrow H^{i}(A^{\bullet}) \longrightarrow \coker f^{i-1}
\longrightarrow \im f^{i} \longrightarrow 0
$$
(These are somehow dual to (1.5.6.3). In fact in some mirror universe this
might have been given as the standard definition of homology.) Assume the
category is that of modules over a fixed ring for convenience, but be aware
that the result is true for any abelian category.
\end{exercise}

\begin{exercise}
EXERCISE AND IMPORTANT DEFINITION.

Suppose
$$
0 \xrightarrow{d^{0}} A^{1} \xrightarrow{d^{1}} \dots \xrightarrow{d^{n-1}}
A^{n} \xrightarrow{d^{n}} 0
$$
is a complex of finite-dimensional k-vector spaces (often called
$A^{\bullet}$ for short). Define $h^i(A^{\bullet}) := \dim H^i(A^{\bullet})$.
Show that $\sum (-1)^i \dim A^i = \sum (-1)^i h^i(A^{\bullet})$. In
particular, if $A^{\bullet}$ is exact, then $\sum (-1)^i \dim A^i = 0$. (If
you haven't dealt much with cohomology, this will give you some practice.)

\begin{proof}
    \begin{align*}
        \sum_{i=1}^{n} \left( -1 \right) ^{i}\dim{A^{i}}
        &= \sum_{i=1}^{n} \left( -1 \right)^{i}\left(
        \dim{\ker{f^{i}}}+\dim{\im{f^{i}}} \right)    \\
        &= \sum_{i=1}^{n} \left( -1 \right) ^{i}\left(
        \dim{\ker{f^{i}}}-\dim{\im{f^{i-1}}} \right) \\
        &= \sum_{i=1}^{n} \left( -1 \right) ^{i}h^{i}\left( A^{\bullet } \right)
    \end{align*}
\end{proof}

\end{exercise}

\begin{exercise}
IMPORTANT EXERCISE.

Suppose $\mathcal{C}$ is an abelian category. Define the category
$Com_{\mathcal{C}}$ of complexes as follows. The objects are infinite
complexes
$$
A^{\bullet}: \quad \dots \longrightarrow A^{i-1} \xrightarrow{f^{i-1}}
A^{i} \xrightarrow{f^{i}} A^{i+1} \xrightarrow{f^{i+1}} \dots
$$
in $\mathcal{C}$, and the morphisms $A^{\bullet} \to B^{\bullet}$ are
commuting diagrams
$$
\begin{tikzcd}
    \cdots \arrow[r] \arrow[d] & A^{i} \arrow[r, "f^{i}"] \arrow[d] & A^{i+1} \arrow[r] \arrow[d] & \cdots \arrow[d] \\
    \cdots \arrow[r] & B^{i} \arrow[r, "g^{i}"] & B^{i+1} \arrow[r] & \cdots
\end{tikzcd}
$$
Show that $Com_{\mathcal{C}}$ is an abelian category.

Feel free to deal with the special case of modules over a fixed ring.
(Remark for experts: Essentially the same argument shows that
$\mathcal{C}^{\mathcal{I}}$ is an abelian category for any small category
$\mathcal{I}$ and any abelian category $\mathcal{C}$. This immediately
implies that the category of presheaves on a topological space $X$ with
values in an abelian category $\mathcal{C}$ is automatically an abelian
category, cf. \S\emph{2.3.5}.)

\begin{proof}
    This is obtained immediately from $\mathscr{C}$ being abelian category.
\end{proof}

\end{exercise}

\begin{exercise}
IMPORTANT EXERCISE.

Show that (1.5.6.5) induces a map of homology $H^{i}(A^{\bullet}) \to
H^{i}(B^{\bullet})$. Show furthermore that $H^{i}$ is a covariant functor
$Com_{C} \to C$. (Again, feel free to deal with the special case
$\mod_{A}$.)

\begin{proof}
    \begin{enumerate}[label=({\roman*})]
        \item 
            By (1.5.6.3) we know any chain complex
            $$
            \ldots \longrightarrow A^{i-1}\xrightarrow[]{f^{i-1}}
            A^{i}\xrightarrow[]{f^{i}}A^{i+1}\xrightarrow[]{f^{i+1}} \ldots 
            $$
            can be decomposed into series of 
            $$
            0\longrightarrow \ker{f^{i}}\longrightarrow A^{i}\longrightarrow
            \im{f^{i}}\longrightarrow 0
            $$
            and
            $$
            0\longrightarrow \im{f^{i-1}}\longrightarrow
            \ker{f^{i}}\longrightarrow H^{i}(A^{\bullet })\longrightarrow 0
            .$$
            Therefore the any homomorphism bewteen chain complexes naturally
            induces homomorphisms such as
            $$
            \ker{f^{i}}\longrightarrow \ker{g^{i}}
            $$
            and
            $$
            \im{f^{i}}\longrightarrow \im{g^{i}}
            .$$
            Hence
            $$
            \ker{f^{i}}/\im{f^{i-1}}\longrightarrow \ker{g^{i}}/\im{g^{i-1}}
            $$
            i.e.
            $$
            H^{i}(A^{\bullet })\longrightarrow H^{i}(B^{\bullet })
            .$$
            Done.
        \item 
            \emph{(I think it's trivial.)}
    \end{enumerate}
\end{proof}

\end{exercise}

\begin{exercise}
Show that two homotopic maps give the same map on homology.

\begin{proof}
    $$
    \begin{tikzcd}
    \ldots \arrow[r] & C^{i-1} \arrow[r, "d_{C}^{i-1}"] \arrow[d, bend right] \arrow[d, bend left] \arrow[ld] & C^{i} \arrow[r, "d_{C}^{i}"] \arrow[d, bend right] \arrow[d, bend left] \arrow[ld] & C^{i+1} \arrow[r, "d_{C}^{i+1}"] \arrow[d, bend right] \arrow[d, bend left] \arrow[ld] & \ldots \arrow[ld] \\
    \ldots \arrow[r] & D^{i-1} \arrow[r, "d_{D}^{i-1}"]                                                       & D^{i} \arrow[r, "d_{D}^{i}"]                                                       & D^{i+1} \arrow[r, "d_{D}^{i+1}"]                                                       & \ldots           
    \end{tikzcd}
    $$
    
    This is euqivalent to prove $f-g$ is gives the zero mapping on homology,
    i.e. $dw-wd$ gives the zero mapping on homology, which is trivial.
\end{proof}

\end{exercise}

\begin{exercise}[Trivial]
Suppose $F$ is an exact functor. Show that applying $F$ to an exact
sequence preserves exactness. For example, if $F$ is covariant, and
$A' \rightarrow A \rightarrow A''$ is exact, then $FA' \rightarrow FA
\rightarrow FA''$ is exact. (This will be generalized in Exercise 1.5.I(c).)

\begin{proof}
    \emph{(Just decompose those sequences)}
\end{proof}

\end{exercise}

\begin{exercise}
Suppose $A$ is a ring, $S \subset A$ is a multiplicative subset, and $M$ is
an $A$-module.

\begin{enumerate}[label=({\alph*})]
\item Show that localization of $A$-modules $\mod_{A} \to \mod_{S^{-1}A}$
    is an exact covariant functor.
\item Show that $(\cdot) \otimes_{A} M$ is a right-exact covariant functor
    $\mod_{A} \to \mod_{A}$. (This is a repeat of Exercise 1.2.H.)
\item Show that $\hom(M,\cdot)$ is a left-exact covariant functor
    $\mod_{A} \to \mod_{A}$. If $\mathcal{C}$ is any abelian category, and
    $C \in \mathcal{C}$, show that $\hom(C,\cdot)$ is a left-exact covariant
    functor $\mathcal{C} \to Ab$.
\item Show that $\hom(\cdot, M)$ is a left-exact contravariant functor
    $\mod_{A} \to \mod_{A}$. If $C$ is any abelian category, and
    $C \in C$, show that $\hom(\cdot, C)$ is a left-exact contravariant
    functor $C \to Ab$.
\end{enumerate}

\begin{proof}
    \begin{enumerate}[label=({\alph*})]
        \item 
            Trivial.
        \item \begin{enumerate}[label=({\roman*})]
            \item 
                See \emph{[Introduction to Commutative Algebra,
                \textbf{Proposition 2.18} -- M. Atiyah]}, trivial.
            \item 
            \end{enumerate}
        \item \begin{enumerate}[label=({\roman*})]
            \item 
                See \emph{[Introduction to Commutative Algebra,
                \textbf{Proposition 2.9} -- M. Atiyah]}, trivial.
            \item 
            \end{enumerate}
        \item \begin{enumerate}[label=({\roman*})]
            \item 
                See \emph{[Introduction to Commutative Algebra,
                \textbf{Proposition 2.9} -- M. Atiyah]}, trivial.
            \item 
            \end{enumerate}
    \end{enumerate}
\end{proof}

\end{exercise}

\begin{exercise}
Suppose $M$ is a finitely presented $A$-module: $M$ has a finite number of
generators, and with these generators it has a finite number of relations;
or usefully equivalently, fits in an exact sequence
$$
A^{\oplus q} \longrightarrow A^{\oplus p} \longrightarrow M \longrightarrow
0.
$$
Use (1.5.9.1) and the left-exactness of Hom to describe an isomorphism
$$
S^{-1}\hom_{A}(M, N) \stackrel{\sim}{\leftrightarrow}
\hom_{S^{-1}A}(S^{-1}M,
S^{-1}N).
$$
(You might be able to interpret this in light of a variant of Exercise
1.5.I below, for left-exact contravariant functors rather than right-exact
covariant functors.)

\begin{proof}
    First, by
    $$
    A^{\bigoplus q}\longrightarrow A^{\bigoplus p}\longrightarrow
    A^{M}\longrightarrow 0  
    $$
    being exact, we have
    $$
    S^{-1}A^{\bigoplus q}\longrightarrow S^{-1}A^{\bigoplus p}\longrightarrow
    S^{-1}A^{M}\longrightarrow 0  
    $$
    being exact, hence
    $$
    0\longrightarrow \hom_{S^{-1}A}{(S^{-1}A^{\bigoplus
    q},S^{-1}N)}\longrightarrow \hom_{S^{-1}A}{(S^{-1}A^{\bigoplus
    p},S^{-1}N)}\longrightarrow \hom_{S^{-1}A}{(S^{-1}M,S^{-1}N)}
    $$
    is exact.

    On the other hand, 
    $$
    A^{\bigoplus q}\longrightarrow A^{\bigoplus p}\longrightarrow
    A^{M}\longrightarrow 0  
    $$
    being exact implies
    $$
    0\longrightarrow \hom_{A}{(A^{\bigoplus q},N)}\longrightarrow
    \hom_{A}{(A^{\bigoplus p}),N}\longrightarrow \hom_{A}{(M,N)}
    $$
    is exact, therefore
    $$ 0\longrightarrow S^{-1}\hom_{A}{(A^{\bigoplus q},N)}\longrightarrow
    S^{-1}\hom_{A}{(A^{\bigoplus p}),N}\longrightarrow S^{-1}\hom_{A}{(M,N)}
    $$
    is exact.

    It's easy to define a natural homomorphism
    $S^{-1}\hom_{A}{(M,N)}\longrightarrow \hom_{S^{-1}A}{(S^{-1}M,S^{-1}N)}$
    explicitly. Also, obviously
    $$
    \hom_{S^{-1}A}{(S^{-1}A,S^{-1}N)}\cong S^{-1}N\text{ by
    }\overline{f}\longmapsto f(1)
    $$
    and
    $$
    S^{-1}\hom_{A}{(M,N)}\cong S^{-1}N\text{ by }\frac{f}{s}\longmapsto
    \frac{f(1)}{s}
    .$$
    Therefore we have the direct sum
    $$
    \hom_{S^{-1}A}{(S^{-1}A^{\bigoplus{r}},S^{-1}N)}\cong 
    \left( S^{-1}N \right) ^{\bigoplus{r}}\cong 
    S^{-1}\hom_{A}{(M^{\bigoplus{r}},N)}
    .$$
\end{proof}

\end{exercise}

\begin{exercise}[Beautiful!]
IMPORTANT EXERCISE (THE FHHF THEOREM).

This result can take you far, and perhaps for that reason it has sometimes
been called the Fernbahnhof (FernbaHnHoF) Theorem, notably in [Vak1, Exer.
1.5.I]. Suppose $F: \mathcal{A} \to \mathcal{B}$ is a covariant functor of
abelian categories, and $C^{\bullet}$ is a complex in $\mathcal{A}$.

\begin{enumerate}[label=({\alph*})]
\item (F right-exact yields $FH^{\bullet} \longrightarrow H^{\bullet}F$)
    If $F$ is right-exact, describe a natural morphism $FH^{\bullet}
    \rightarrow H^{\bullet}F$. (More precisely, for each $i$, the left side
    is $F$ applied to the cohomology at piece $i$ of $C^{\bullet}$, while
    the right side is the cohomology at piece $i$ of $FC^{\bullet}$.)
\item (F left-exact yields $FH^{\bullet} \leftarrow H^{\bullet}F$) If $F$
    is left-exact, describe a natural morphism $H^{\bullet}F \rightarrow
    FH^{\bullet}$.
\item (F exact yields $FH^{\bullet} \xleftarrow{\sim} H^{\bullet}F$) If $F$
    is exact, show that the morphisms of (a) and (b) are inverses and thus
    isomorphisms.
\end{enumerate}

Hint for (a): use $C^i \xrightarrow{d^i} C^{i+1} \longrightarrow
\text{coker } d^i \longrightarrow 0$ to give an isomorphism $F \text{coker }
d^i \stackrel{\sim}{\leftrightarrow} \text{coker } F d^i$. Then use the first line of
(1.5.6.4) to give an epimorphism $F \text{im } d^i \twoheadrightarrow
\text{im } F d^i$. Then use the second line of (1.5.6.4) to give the
desired map $FH^i C^\bullet \to H^i FC^\bullet$. While you are at it, you
may as well describe a map for the fourth member of the quartet {coker, im,
H, ker}: $F \text{ker } d^i \to \text{ker } F d^i$.

\begin{proof}
    \emph{(Note that in homological algebra, it's a consequence that every left,
    right exact, or exact, functor being additive. Therefore it's well-defined
    to write} $H^{\bullet}F$\emph{.)}

    \begin{enumerate}[label=({\roman*})]
        \item 
            We follow the hint: consider the exact sequence
            $$
            C^{i}\xrightarrow[]{d^{i}} C^{i+1}\longrightarrow
            \coker{d^{i}}\longrightarrow 0
            $$
            and apply it with functor $F$, we have
            $$
            F\left( C^{i} \right) \xrightarrow[]{F\left( d^{i} \right) } F\left(
            C^{i+1} \right) \longrightarrow F\coker{d^{i}}\longrightarrow 0
            $$
            being exact. But at the same time we have the exact sequence
            $$
            F\left( C^{i} \right) \xrightarrow[]{F\left( d^{i} \right) } F\left(
            C^{i+1} \right) \longrightarrow \coker{F{d^{i}}}\longrightarrow 0
            .$$
            Hence $F{\coker{d^{i}}}\cong \coker{F{d^{i}}}$.

            From (1.5.6.4), we know
            $$
            0\longrightarrow \im{d^{i}}\longrightarrow C^{i+1}\longrightarrow
            \coker{d^{i}}\longrightarrow 0
            $$
            is exact. Apply it with $F$,
            $$
            F{\im{d^{i}}}\longrightarrow F{C^{i+1}}\longrightarrow
            F{\coker{d^{i}}}\longrightarrow 0
            $$
            is exact.
            $$
            \begin{tikzcd}
                                                        &                          &                                 & \coker{F{d^{i}}} \arrow[rd]                           &   \\
            F{\im{d^{i}}} \arrow[rr] \arrow[rd, dotted] &                          & F{C^{i+1}} \arrow[r] \arrow[ru] & F{\coker{d^{i}}} \arrow[u, no head, dashed] \arrow[r] & 0 \\
                                                        & \im{F{d^{i}}} \arrow[ru] &                                 &                                                       &  
            \end{tikzcd}
            $$
            and we have the epimorphism
            $$
            F{\im{d^{i}}}\longrightarrow \im{F{d^{i}}}\longrightarrow 0
            $$

            Again from (1.5.6.4), we know
            $$
            0\longrightarrow H^{i}\left( C^{\bullet} \right) \longrightarrow
            \coker{d^{i-1}}\longrightarrow \im{d^{i}}\longrightarrow 0
            $$
            is exact. Apply it with F,
            $$
            \begin{tikzcd}
                        & F{H^{i}{(C^{\bullet})} } \arrow[r] \arrow[d, dotted] & F\left( \coker{d^{i-1}} \right) \arrow[r] \arrow[d, "\sim"] & F\left( \im{d^{i}} \right) \arrow[r] \arrow[d, two heads] & 0 \\
            0 \arrow[r] & H^{i}F{(C^{\bullet})} \arrow[r]                      & \coker{F{d^{i-1}}} \arrow[r]                                & \im{F{d^{i}}} \arrow[r]                                   & 0
            \end{tikzcd}
            $$
            we obtain the natural morphism by universal property.
        \item 
            Similar to \emph{(i)}, we use the two short exact sequences in
            (1.5.6.3) and obtain
            $$
            \begin{tikzcd}
            0 \arrow[r] & \im{F{d^{i-1}}} \arrow[r] \arrow[d, hook] & \ker{F{d^{i}}} \arrow[r] \arrow[d, "\sim"] & H^{i}{(F{C^{\bullet}})} \arrow[r] \arrow[d, dotted] & 0 \\
            0 \arrow[r] & F{\im{d^{i-1}}} \arrow[r]                 & F{\ker{d^{i}}} \arrow[r]                   & F{H^{i}{(C^{\bullet})}}                             &  
            \end{tikzcd}
            .$$
            Done
        \item 
            If $F$ is exact, the epimorphism
            $$
            F{\im{d^{i}}}\longrightarrow \im{F{d^{i}}}\longrightarrow 0
            $$
            obtained in \emph{(i)} becomes isomorphism
            $$
            F{\im{d^{i}}}\cong\im{F{d^{i}}}
            ,$$
            therefore the induced natural morphism is an isomorphism. The same
            story goes to \emph{(ii)}. Thus the morphism in \emph{(i)} and
            \emph{(ii)}, denoted by $\alpha,\beta$, are both isomorphisms. Now
            we'll prove they're acturally inverses to each other:

            Consider
            $$
            \begin{tikzcd}
            0 \arrow[r] & \ker{F(d^i)} \arrow[r, two heads] \arrow[d, "\sim"]  & H^iF(C^\bullet) \arrow[r, hook] \arrow[d, "\beta"]    & \coker{F(d^{i-1})} \arrow[d, "\sim"] \arrow[r] & 0 \\
            0 \arrow[r] & F(\ker{d^i}) \arrow[r, two heads]                    & F(H^i(C^\bullet)) \arrow[r, hook] \arrow[d, "\alpha"] & F(\coker{d^{i-1}}) \arrow[d, "\sim"] \arrow[r] & 0 \\
            0 \arrow[r] & \ker{F(d^i)} \arrow[u, "\sim"'] \arrow[r, two heads] & H^iF(C^\bullet) \arrow[r, hook]                       & \coker{F(d^{i-1})} \arrow[r]                   & 0
            \end{tikzcd}
            .$$
            We have
            $$
            (\text{mono}) \circ (\alpha \circ \beta) \circ (\text{epi}) =
            (\text{mono}) \circ \text{id}_{H^iF(C^\bullet)} \circ (\text{epi})
            $$
            By the cancellation law of monomorphism and epimorphism, we have
            $$
            \alpha\circ \beta=\id
            $$
    \end{enumerate}
\end{proof}

\end{exercise}

\begin{exercise}
EXERCISE (Kernels COMMUTE WITH LIMITS).

Suppose $\mathcal{C}$ is an abelian category, and $a\colon \mathcal{I} \to
\mathcal{C}$ and $b\colon \mathcal{I} \to \mathcal{C}$ are two diagrams in
$\mathcal{C}$ indexed by $\mathcal{I}$. For convenience, let $A_i = a(i)$
and $B_i = b(i)$ be the objects in those two diagrams. Let $h_i\colon A_i
\to B_i$ be maps commuting with the maps in the diagram. (Translation: $h$
is a natural transformation of functors $a \to b$, see \S\emph{1.1.21}.)
Then the $\ker h_i$ form another diagram in $\mathcal{C}$ indexed by
$\mathcal{I}$. Describe a canonical isomorphism $\lim \ker h_i
\xleftarrow{\sim} \ker(\lim A_i \to \lim B_i)$, assuming the limits exist.

\begin{proof}
    Chase the diagrams and obtain those commute diagrams:
    $$
    \begin{tikzcd}
                                                                                                              & i \arrow[d] \arrow[rd]                                                               &                     \\
    \ker{h_i} \arrow[r]                                                                                       & A_i \arrow[r]                                                                        & B_i                 \\
    \lim{\ker{h_i}} \arrow[u] \arrow[r] \arrow[d, "\exists!", dotted] \arrow[rr, "0" description, bend right] & \lim{A_i} \arrow[r, "\exists!"] \arrow[u] \arrow[lu, "\exists!" description, dotted] & \lim{B_i} \arrow[u] \\
    \ker{(\lim{A_i}\rightarrow\lim{B_i})} \arrow[ru]                                                            &                                                                                      &                    
    \end{tikzcd}
    $$
    $$
    \begin{tikzcd}
                                                                           & C \arrow[d] \arrow[rd] \arrow[rrd, "0"]                       &                         &               \\
    \lim{A_i} \arrow[r, dotted] \arrow[rd, "\exists!" description, dotted] & \ker{h_i} \arrow[r] \arrow[d, "\exists!" description, dotted] & A_i \arrow[r] \arrow[d] & B_i \arrow[d] \\
                                                                           & \ker{h_j} \arrow[r]                                           & A_j \arrow[r]           & B_i           \\
    \ker{(\lim{A_i}\rightarrow\lim{B_i})} \arrow[uu]                       &                                                               &                         &              
    \end{tikzcd}
    $$
\end{proof}

\end{exercise}

\begin{exercise}
Make sense of the statement that ``limits commute with limits'' in a
general category, and prove it. (Hint: recall that kernels are limits. The
previous exercise should be a corollary of this one.)
\end{exercise}

\end{document}
